\documentclass{article}

\usepackage[utf8]{inputenc}

\title{Compte rendu de la réunion du 24 Avril}
\date{}

\begin{document}
	
\maketitle
	
\section*{Premiers résultats}
	\begin{itemize}
		\item Les attaques statiques sont meilleurs que les dynamiques
		\begin{itemize}
			\item[$\rightarrow$] Hypothèses de départ peuvent êtres fausses
			\item[$\rightarrow$] Soucis liés à la simulation (implémentation, trop éloignée de la réalité)
		\end{itemize}
		\item Pour la suite des simus
		\begin{itemize}
			\item Refaire des simulations en faisant varier le delta et comparer au random statique
			\item Modifier l'algo pour qu'un lien soit traversable uniquement s'il n'est pas bloqué durant toute la durée de la traversée
			\item Regarder l'accessibilité pour avoir un autre point de vue
			\item Essayer de faire des heatmap sur des petites villes (à cause du nombre de simulations nécessaires)
		\end{itemize}		
	\end{itemize}
	
\section*{Optimisations}
	\begin{itemize}
		\item On a optimisé Dijkstra pour qu'il écrive plus de résultats
		\item On a une invalidation de liens en cours de développement
	\end{itemize}

% \newpage
	
\section*{Pour la suite des simus}
	\begin{itemize}
		\item Refaire des simulations en faisant varier le delta
		\item Modifier l'algo pour qu'un lien soit traversable uniquement s'il n'est pas bloqué durant toute la durée de la traversée
		\item Essayer de faire des heatmap sur des petites villes (à cause du nombre de simulations nécessaires)
	\end{itemize}

\section*{Compte-rendu}
	\begin{itemize}
		\item Intro : Contexte pour présenter le sujet et problématique
		\item Notions de bases (rappels sur les graphes)
		\item Trois parties pour montrer la répartition \begin{itemize}
			\item Traitement du graphe
			\item Génération des attaques
			\item Mesures de l'efficacité
		\end{itemize}
		\item Perspectives réservé à ce que l'on a pas pu faire
		
		$\rightarrow$ Ne pas sortir un nouvel algo
		
		\item Les figures doivent êtres accompagnées d'une description et d'un commentaire, et doivent servir le propos
		\item Imprimer le compte-rendu pour voir ce que ça donne
	\end{itemize}

\section*{Mise en forme des plots}
	\begin{itemize}
		\item Pas plus de 5 courbes par figures
		\item Toute les fig doivent avoir une description et un commentaire
		\item Légende, axes cohérent (noms également)
		\item On peut également s'arranger pour que les courbes ait des allures différentes
		\item Essayer différents types de plot
		
		$\rightarrow$ Par exemple violon plot pour les distributions aléatoires
		\item Figures actuelles : \begin{itemize}
			\item Titres (et légendes) beaucoup trop petits
			\item Trop d'attaques
		\end{itemize}
		\item On peut aussi faire des plots pour mesurer le temps de calcul
		
		Question typique : "combien de temps ça prend pour ma  ville ?"
	\end{itemize}

\section*{Oral}
	\begin{itemize}
		\item 15 minutes
		\item Diapo : \begin{itemize}
			\item Pas plus d'une diapo par minute
			\item Pas de paragraphes
		\end{itemize}
		\item Essayer de répartir le temps et de laisser chacun parler de sa partie
		\item On présente le principe du nouvel algo (mais pas l'ancien)
		\item Conclusion : \begin{itemize}
			\item Dernière diapo (pas de diapo "merci d'avoir écouté")
			\item Phrases/points sur lesquelles le jury devras rebondir
		\end{itemize}
	\end{itemize}
\end{document}
